\begin{abstract}
Modern intelligent systems increasingly rely on multiple components, including large language models (LLMs), retrieval systems, workflow engines, knowledge stores, and user feedback mechanisms. While each component has advanced significantly, they are often developed as independent subsystems with limited coordination. Workflow engines optimize execution, retrieval systems manage information access, and feedback mechanisms are frequently confined to isolated evaluation pipelines. As a result, continual adaptation across the entire system remains fragmented.
This paper introduces the \emph{Semut Adaptive Architecture}, a unified systems architecture that models intelligent systems through three interconnected memory structures: \emph{Artifact Memory}, \emph{Action Graph}, and \emph{Feedback Memory}. Artifact Memory preserves persistent contextual information such as documents, source code, conversations, and structured knowledge. Action Graph represents executable actions and their dependencies as an evolving graph of operational behavior. Feedback Memory continuously records human evaluations, execution outcomes, environmental observations, and system performance, enabling future decisions to incorporate accumulated experience.
Rather than replacing existing reasoning systems, the proposed architecture is reasoning-engine independent. Rule-based systems, large language models, future learning algorithms, and human operators can all utilize the same adaptive memory framework. Retrieved artifacts, historical actions, and prior feedback collectively provide contextual evidence for selecting subsequent actions. The outcomes of execution are subsequently incorporated back into all three memory structures, forming a continuous adaptive loop.
The primary contribution of this work is not a new reasoning algorithm, but a unified architectural abstraction that integrates artifacts, actions, and feedback into a coherent adaptive memory system. This abstraction provides a common interface for continual learning, workflow execution, retrieval, and human-in-the-loop collaboration across heterogeneous intelligent systems. Potential applications include personal AI assistants, software engineering agents, enterprise automation, robotics, and multi-agent coordination.
\end{abstract}